\documentclass[aps,pra,twocolumn,reprint,superscriptaddress,nofootinbib]{revtex4-2}

  \usepackage[svgnames,table,xcdraw]{xcolor}
  \usepackage{amsmath}
  \usepackage{microtype}
  \usepackage{amssymb}
  \usepackage{braket}
  \usepackage{graphicx}
  \usepackage{booktabs}
%\usepackage{tabu}
  \usepackage{bm}
  \usepackage{dsfont}
  \usepackage{hyperref}
  \usepackage[caption=false]{subfig}
  \usepackage{amsthm}
  \newtheorem*{claim}{Claim}
  \usepackage{xfrac}
  \usepackage{enumitem}
  \usepackage{orcidlink}
  \usepackage{appendix}
  \usepackage[ruled,vlined]{algorithm2e}

%\usepackage{array}
%\newcolumntype{L}[1]{>{\raggedright\arraybackslash}p{#1}}

% \usepackage[
%     backend=biber,
%     style=numeric-comp,
%     maxbibnames=2,
%     sorting=none,
%     url=false
% ]{biblatex}
% \addbibresource{refs.bib}
% \ExecuteBibliographyOptions{doi=false}
% \newbibmacro{string+doi}[1]{%
%   \iffieldundef{doi}{#1}{\href{http://dx.doi.org/\thefield{doi}}{#1}}}
% \DeclareFieldFormat{title}{\usebibmacro{string+doi}{\mkbibemph{#1}}}
% \DeclareFieldFormat[article]{title}{\usebibmacro{string+doi}{\mkbibquote{#1}}}


  \hypersetup{colorlinks,linkcolor={DarkBlue},citecolor={DarkBlue},urlcolor={DarkBlue}}

  \let\originalleft\left
                     \let\originalright\right
  \renewcommand{\left}{\mathopen{}\mathclose\bgroup\originalleft}
  \renewcommand{\right}{\aftergroup\egroup\originalright}

  \usepackage{mathtools}
  \DeclarePairedDelimiter{\paren}{(}{)}

  \newcommand{\id}{\mathds{1}}
  \DontPrintSemicolon
  \SetKwInput{KwInput}{Input}
  \SetKwInput{KwOutput}{Output}
  \SetKwFor{ForEach}{foreach}{do}{}
  \SetKwFor{ForLoop}{for}{do}{}

  % !TEX root =main.tex
\usepackage{nag}
\usepackage{graphicx}
\usepackage{amsthm}
% \usepackage{tabulary} % Unused; can interfere with standard tabular in RevTeX
\usepackage{qcircuit}
\usepackage{mathtools}
\usepackage{bm}
\usepackage[ruled,vlined]{algorithm2e}
%\usepackage[table,xcdraw]{xcolor}
\usepackage{braket}
\usepackage{datetime}
\usepackage{stackrel}
% Used to declare symbols below
% \usepackage{stmaryrd}
\usepackage{dsfont}
\usepackage{qcircuit}
\usepackage[english]{babel}
\usepackage{units}

\usepackage{tikz}
\usetikzlibrary{backgrounds,decorations.pathreplacing}

\usepackage{chngcntr}
\counterwithout{equation}{section}

\usepackage[normalem]{ulem}
\usepackage{slashed}
%%
\usepackage{soul}

\usepackage{xcolor}
% \usepackage[
%     colorlinks=true,
%     urlcolor=blue,
%     linkcolor=blue,
%     citecolor=blue,
%     filecolor=blue,
% ]{hyperref}
% \usepackage[caption=false]{subfig}
%%
\usepackage{amssymb}
\usepackage{amsmath}
\usepackage{amsthm}
\usepackage{braket}
\usepackage{mathdots}

% \usepackage[OT2,T1]{fontenc}
% \DeclareSymbolFont{cyrletters}{OT2}{wncyr}{m}{n}
% \DeclareMathSymbol{\Sha}{\mathalpha}{cyrletters}{"58}






\usepackage{makeidx}
\makeindex


% 

\usepackage{soul}
\usepackage{xargs}
%\usepackage{todonotes}
\newcommandx{\cmnote}[2][1=]{\linespread{1.0}\todo[linecolor=red,backgroundcolor=red!25,bordercolor=red,#1]{#2}}

\newcommand{\CS}{\text{CS}}
\newcommand{\vac}{\text{vac}}

% 
\let\underline\ul

% 

\allowdisplaybreaks[4]


\DeclareMathOperator{\tr}{tr}
\DeclareMathOperator{\Tr}{Tr}
\DeclareMathOperator{\TR}{TR}
\DeclareMathOperator{\Sp}{Sp}
\DeclareMathOperator{\spur}{sp}
\DeclareMathOperator{\diag}{diag}
\DeclareMathOperator{\Herm}{Herm}
\DeclareMathOperator{\cov}{cov}
\DeclareMathOperator{\var}{var}
\DeclareMathOperator{\Prob}{Prob}
\DeclareMathOperator{\spec}{spec}
\DeclareMathOperator{\sinc}{sinc}
\renewcommand{\Re}{\operatorname{Re}}
\renewcommand{\Im}{\operatorname{Im}}

\DeclareMathOperator{\sech}{sech}
\DeclareMathOperator{\csch}{csch}

\newcommand{\diffd}{\mathrm{d}}



 \index{} \index{} \index{} \index{} \index{} \index{} \index{} \index{}%%
\makeatletter

%

%
\newcommand{\ringplus}{\mathbin{\text{\@ringplus}}}

\newcommand{\@ringplus}{%
  \ooalign{\hidewidth\raise1.3ex\hbox{\tiny$\circ$}\hidewidth\cr$\m@th+$\cr}%
}

\newcommand{\ringminus}{\mathbin{\text{\@ringminus}}}

\newcommand{\@ringminus}{%
  \ooalign{\hidewidth\raise0.9ex\hbox{\tiny$\circ$}\hidewidth\cr$\m@th-$\cr}%
}
\makeatother
 \index{} \index{} \index{} \index{} \index{} \index{} \index{} \index{}%%

\newcommand{\rplus}{{\ringplus}}
\newcommand{\rminus}{{\ringminus}}

\newcommand{\nasc}{\epsilon}
\newcommand{\dnasc}{\delta_\nasc}
\newcommand{\tdnasc}{\tilde \delta_\nasc}

\newcommand{\qnasc}{q_\nasc}
\newcommand{\pnasc}{p_\nasc}


\newcommand{\mc}[1]{\mathcal{#1}}
\newcommand{\const}[0]{\text{(const.)}}
\newcommand{\cc}[0]{\text{c.c.}}
\newcommand{\bvec}[1]{\bm{#1}}

\newcommand{\tp}[0]{T}
\newcommand{\herm}[0]{\mathrm{H}}
\newcommand{\psinv}[0]{+}

\newcommand{\group}[1]{\mathrm{{#1}}}
\newcommand{\algebra}[1]{\mathfrak{{#1}}}
\newcommand{\logical}[0]{L}

\newcommand{\mathlatt}[1]{{\mathsf{#1}}}
\newcommand{\mathreclatt}[1]{\tilde{\mathsf{#1}}}
\newcommand{\dirlatt}{{\mathlatt L}}
\newcommand{\reclatt}{{\tilde \dirlatt}}
\newcommand{\vlatt}{{\mathlatt V}}
\newcommand{\flatt}{{\mathlatt F}}
\newcommand{\slatt}{{\mathlatt S}}


\newcommand{\weyl}{V}

\newcommand{\EPR}{\text{EPR}}
\newcommand{\bounceEPR}[2]{\op{A}_{#2}(#1)}
\newcommand{\bounceEPRgate}[2]{A_{#2}(#1)}

\newcommand{\dampop}[1]{\op{N}(#1)}
\newcommand{\dampgate}{N(\beta)}

% GKP commands
\newcommand{\GKP}{\text{GKP}} 
\newcommand{\GKPproj}{\op{\Pi}_\GKP}
\newcommand{\approxGKPproj}{\hat{\overline{\Pi}}_\GKP}
\newcommand{\qunaught}{\varnothing}
\newcommand{\sqz}{\grn \text{sqz}\blk}


\newcommand{\mvec}[1]{\tilde{\vec{#1}}}

\newcommand{\latG}{V}
%---------------------------------------------------------------------------
% Author commenting
%---------------------------------------------------------------------------




%---------------------------------------------------------------------------
% Sha character (Dirac comb)
%---------------------------------------------------------------------------

\DeclareFontFamily{U}{wncy}{}
\DeclareFontShape{U}{wncy}{m}{n}{<->wncyr10}{}
\DeclareSymbolFont{mcy}{U}{wncy}{m}{n}
\DeclareMathSymbol{\Sh}{\mathord}{mcy}{"58}



%---------------------------------------------------------------------------
% Specialized bra and ket defs
%---------------------------------------------------------------------------

% NOTE: basic bras and kets already defined in braket

% negative spacing for subscripts -- decide on one or two \! characters:
\newcommand{\negspace}{\!}
%\newcommand{\negspace}[0]{\!\!}

% Base commands for putting left and right subscripts on bras and kets and things
\newcommand{\lsub}[2]{{\protect\vphantom{#1}}_{#2} \negspace {#1}}
\newcommand{\rsub}[2]{{#1} \negspace {\protect\vphantom{#1}}_{#2}}
\newcommand{\lrsub}[3]{{\protect\vphantom{#1}}_{#2} \negspace {#1} \negspace {\protect\vphantom{#1}}_{#3}}

% Subscripted bras/kets
\newcommand{\ketsub}[2]{\rsub {\ket{#1}} {#2}}
\newcommand{\brasub}[2]{\lsub {\bra{#1}} {#2}}

% bra/ket macros
\newcommand{\pbra}[1]{\brasub{#1} p}
\newcommand{\qbra}[1]{\brasub{#1} q}
\newcommand{\pket}[1]{\ketsub{#1} p}
\newcommand{\qket}[1]{\ketsub{#1} q}



% basic inner product
\newcommand{\inprod}[2]{\left\langle {#1} | {#2} \right\rangle}
%\renewcommand{\braket}[0]{\inprod}
\newcommand{\mathinprod}[2]{\avg{{#1},{#2}}}

% doubly subscripted inner product
\newcommand{\inprodsubsub}[4]{\lrsub {\inprod{#1}{#2}} {#3} {#4}}

% inner product macros
\newcommand{\pqinprod}[2]{\inprodsubsub{#1}{#2} p q}
\newcommand{\qpinprod}[2]{\inprodsubsub{#1}{#2} q p}
\newcommand{\pinprod}[2]{\inprodsubsub{#1}{#2} p p}
\newcommand{\qinprod}[2]{\inprodsubsub{#1}{#2} q q}

%basic outer product
\newcommand{\outprod}[2]{\ket {#1}\!\bra {#2}}

% doubly subscripted outer product
\newcommand{\outprodsubsub}[4]{\ketsub {#1}{#3} \brasub{#2}{#4}}

% outer product macros
\newcommand{\pqoutprod}[2]{\outprodsubsub{#1}{#2}p q}
\newcommand{\qpoutprod}[2]{\outprodsubsub{#1}{#2}q p}
\newcommand{\qoutprod}[2]{\outprodsubsub{#1}{#2}q q}
\newcommand{\poutprod}[2]{\outprodsubsub{#1}{#2}p p}


%---------------------------------------------------------------------------
% Nascent-delta-based bra and ket defs
%---------------------------------------------------------------------------


% bra/ket macros
\newcommand{\pnascbra}[1]{\brasub{#1} \pnasc}
\newcommand{\qnascbra}[1]{\brasub{#1} \qnasc}
\newcommand{\pnascket}[1]{\ketsub{#1} \pnasc}
\newcommand{\qnascket}[1]{\ketsub{#1} \qnasc}

% inner product macros
\newcommand{\pqnascinprod}[2]{\inprodsubsub{#1}{#2} \pnasc \qnasc}
\newcommand{\qpnascinprod}[2]{\inprodsubsub{#1}{#2} \qnasc \pnasc}
\newcommand{\pnascinprod}[2]{\inprodsubsub{#1}{#2} \pnasc \pnasc}
\newcommand{\qnascinprod}[2]{\inprodsubsub{#1}{#2} \qnasc \qnasc}

% outer product macros
\newcommand{\pqnascoutprod}[2]{\outprodsubsub{#1}{#2}\pnasc \qnasc}
\newcommand{\qpnascoutprod}[2]{\outprodsubsub{#1}{#2}\qnasc \pnasc}
\newcommand{\qnascoutprod}[2]{\outprodsubsub{#1}{#2}\qnasc \qnasc}
\newcommand{\pnascoutprod}[2]{\outprodsubsub{#1}{#2}\pnasc \pnasc}




%---------------------------------------------------------------------------
% end of subscripted bra/ket defs
%---------------------------------------------------------------------------

\newcommand{\opbra}[1]{\left( {#1} \right\rvert}
\newcommand{\opket}[1]{\left\lvert {#1} \right)}
\newcommand{\opinprod}[2]{\left( {#1} \middle\vert {#2} \right)}
\newcommand{\opoutprod}[2]{\opket {#1} \opbra {#2}}



\newcommand{\avg}[1]{\left\langle {#1} \right\rangle}
\newcommand{\avgg}[1]{\langle {#1} \rangle}
\newcommand{\abs}[1]{\left\lvert{#1}\right\rvert}
\newcommand{\abss}[1]{\lvert{#1}\rvert}
\newcommand{\norm}[1]{\left\lVert{#1}\right\rVert}
\newcommand{\normm}[1]{\lVert{#1}\rVert}



\newcommand{\fpart}[1]{\left\{{#1}\right\}}
\newcommand{\fpartt}[1]{\{{#1}\}}

%\newcommand{\ipart}[1]{\left\lfloor{#1}\right\rfloor}
%\newcommand{\ipartt}[1]{\lfloor{#1}\rfloor}

\newcommand{\ipart}[1]{\left\lfloor{#1}\right\rceil}
\newcommand{\ipartt}[1]{\lfloor{#1}\rceil}


\newcommand{\cent}{C}
\newcommand{\sat}{S}



\DeclareMathOperator{\pv}{p.v.}




\newcommand{\reals}[0]{\mathbb{R}}
\newcommand{\complex}[0]{\mathbb{C}}
\newcommand{\wholenums}[0]{\mathbb{N}_0}
\newcommand{\integers}[0]{\mathbb{Z}}
\newcommand{\rationals}[0]{\mathbb{Q}}
\newcommand{\projective}{\mathbb{P}}


\newcommand{\op}[1]{\hat{#1}}
\newcommand{\ope}[1]{\hat{#1}}
%\newcommand{\op}[1]{#1}
\newcommand{\opvec}[1]{\op{\vec{#1}}}
\newcommand{\opmat}[1]{\op{\mat{#1}}}
%\newcommand{\id}[0]{1}
%\newcommand{\id}[0]{I}
\newcommand{\supop}[1]{\mathcal{#1}}

%\newcommand{\1}[0]{\mathit{1}}

%\newcommand{\1}[0]{\mat{\id}}

%\newcommand{\mat}[1]{\mathsf{#1}}
%\newcommand{\mat}[1]{\vec{#1}}
\newcommand{\mat}[1]{\bm{\mathrm{#1}}}
\newcommand{\barmat}[1]{\bar{\mat{#1}}}
\newcommand{\tildemat}[1]{\tilde{\mat{#1}}}
%\newcommand{\mat}[1]{{#1}}
\newcommand{\barvec}[1]{\bar{\vec{#1}}}
\newcommand{\tildevec}[1]{\tilde{\vec{#1}}}

%\renewcommand{\vec}[1]{\mathrm{#1}}
\renewcommand{\vec}[1]{\bm{#1}}

\newcommand{\controlled}[1]{\op{\mathrm{C}}_{#1}}
\newcommand{\CZ}[0]{\controlled Z}
\newcommand{\CX}[0]{\controlled X}
\newcommand{\CZct}[0]{\controlled Z^\rightarrow}
\newcommand{\CXct}[0]{\controlled X^\rightarrow}
\newcommand{\CZtc}[0]{\controlled Z^\leftarrow}
\newcommand{\CXtc}[0]{\controlled X^\leftarrow}

\newcommand{\opCZ}[0]{\op{\mathrm{C}}_{Z}}

\newcommand{\cH}[0]{\mathcal{H}}






 \index{} \index{} \index{} \index{} \index{} \index{} \index{} \index{}%%
% COLOUR COMMANDS
 \index{} \index{} \index{} \index{} \index{} \index{} \index{} \index{}%%
\newcommand{\grn}{\color{green!60!black}}
\newcommand{\red}{\color{red!90!black}}
\newcommand{\blu}{\color{blue!70!cyan}}
\newcommand{\blk}{\color{black}}
\newcommand{\gra}{\color{gray!50!white}}
\newcommand{\prp}{\color{purple!70!black}}
\newcommand{\org}{\color{orange!80!black}}
\newcommand{\pnk}{\color{pink}}

 \index{} \index{} \index{} \index{} \index{} \index{} \index{} \index{}%%
% VERY SPECIFIC COMMANDS
 \index{} \index{} \index{} \index{} \index{} \index{} \index{} \index{}%%
\usepackage{graphicx}
%\graphicspath{{./figures/}}

%Circuit macros
\newcommand{\bs}{\! B_{_{\! \downarrow}} \!\!}
%\newcommand{\bsop}{\op{B}_{_{\! \downarrow}}}
\newcommand{\bsop}{\op{B} }
% Figures Location

\newcommand{\rket}[1][0]{\ketsub{#1}{\theta}}
\newcommand{\rbra}[1][0]{\brasub{#1}{\theta}}
\usepackage{graphicx}
\usepackage{multirow}
\usepackage{hhline}
%\graphicspath{{../Figures/}}



%%%%%%%%%%%%%%%%%%%%%%%%%%%%%%%
% beamsplitter as arrow
% Use \bsbal{} just like \ctrl{}
%%%%%%%%%%%%%%%%%%%%%%%%%%%%%%%
\xyoption{color}
\xyoption{line}

\newcommand{\qwax}[1][-1]{\ar @{->} [#1,0]}
%Using the bsbal command requires setting the length. Specifying an optional argument before the length allows to change the colour, specifying an optional argument AFTER the length allows to change the line style (eg. .>, -->)
\newcommandx*\bsbal[3][1=black, 3=->]{\ar @[#1]@{#3} [#2,0] \qw}

\newcommandx*\varbs[5][1=black, 3=\theta,4=0.5,5=->]{\ar @[#1]@{#5}^(#4){#3} [#2,0] \qw}

%\newcommand{\bsbal}[1]{\qwax[{#1}] \qw}
\newcommand{\qwaxred}[1][-1]{\ar @[red]@{->} [#1,0]}
%\newcommand{\bsbalred}[1]{\qwaxred[{#1}] \qw}
\newcommand{\qwred}[1][-1]{\ar @[red]@{-} [0,#1]}
\newcommand{\gatered}[1]{*+<.6em>{#1} \POS ="i","i"+UR;"i"+UL **\dir{-};"i"+DL **\dir{-};"i"+DR **\dir{-};"i"+UR **\dir{-},"i" \qwred}


\newcommand{\ctrlnowire}[1]{\control \qwx[#1]}
\newcommandx*\lblline[3][3=0.5]{\ar @{-}^(#3){#1} [#2,0]}
% \newcommandx*\ctrlg[2]{\control \ar @{-}^{#1} [#2,0] \qw}
\newcommandx*\ctrlg[3][3=0.5]{ \raisebox{-3pt}{$\bullet$}  \ar @{-}^(#3){#1} [#2,0] \qw }
\newcommandx*\ctrlog[2]{\controlo \ar @{-}^{#1} [#2,0] \qw}
%%%Controlo gate followed by dashed wire
\newcommandx*\ctrlodash[1]{\controlo \ar @{-} [#1,0] \ar @[black]@{.} [0,-1]}
\newcommand{\qwdot}{\ar @[black]@{.} [0,-1]}
%%%%%%%%%%%%%%%%%%%%%%%%%%%%%%%

%%%%%%%%%%%%%%%%%%%%%%%%%%%%%%%
%EPR state angled wire
%use \EPR followed by u, d for up and down and then r, l for right and left
%%%%%%%%%%%%%%%%%%%%%%%%%%%%%%%

\newcommand{\EPRdr}{\ar @{-} [dr(0.5)]\qw }
\newcommand{\EPRur}{\ar @{-} [ur(0.5)]\qw }
\newcommand{\EPRul}{\ar @{-} [ul(0.5)] }
\newcommand{\EPRdl}{\ar @{-} [dl(0.5)] }

%%%%%%%%%% invisible gate box %%%%%%%%%%%%%
\newcommand{\nogate}[1]{*+<.6em>{#1} \POS ="i","i"+UR;"i"+UL **[white]\dir{-};"i"+DL **[white]\dir{-};"i"+DR **[white]\dir{-};"i"+UR **[white]\dir{-},"i" }

%%%%%%%%%%%%% Multigate invisible, use \nghost for lower sections, not \ghost %%%%%%%%%%%%%%
\newcommand{\nomultigate}[2]{*+<1em,.9em>{\hphantom{#2}} \POS [0,0]="i",[0,0].[#1,0]="e",!C *{#2},"e"+UR;"e"+UL **[white]\dir{-};"e"+DL **[white]\dir{-};"e"+DR **[white]\dir{-};"e"+UR **[white]\dir{-},"i"}

\newcommand{\colorgategroup}[7]{\POS"#1,#2"."#3,#2"."#1,#4"."#3,#4"!C*+<#5>[#7]\frm{#6}}

\setcounter{page}{1}
\pagenumbering{arabic}

\newcommand{\prlsection}[1]{\vspace{0.5ex}\textbf{#1}}

\newcommand*\circled[1]{\tikz[baseline=(char.base)]{
            \node[shape=circle,draw,inner sep=2pt] (char) {#1};}}


  \begin{document}

    \title{%
      \texorpdfstring
      {SBUP$^3$: A Physics-Calibrated Time-Dependent Density Functional Theory--Semiconductor Bloch Equation Framework for Ultrafast \\Laser--Matter Dynamics in Gallium Arsenide (GaAs)}
      {SBUP3: A Physics-Calibrated Time-Dependent Density Functional Theory-Semiconductor Bloch Equation Framework for Ultrafast Laser-Matter Dynamics in Gallium Arsenide (GaAs)}
    }
    \def \affGatech {College of Computing, Georgia Institute of Technology, Atlanta, GA 30332 USA}
    \def \fora {Department of Computer Science, Federal University of Juiz de Fora, Juiz de Fora MG 36036-900, Brazil}
    \def\ganj{Department of BioNano Technology, Gachon University, 1342 Seongnamdae-ro, Sujeong-gu, Seongnam, Gyeonggi, 13120, Republic of Korea}


    \author{Dennis Delali Kwesi Wayo \orcidlink{0000-0001-9980-6247}}
    \affiliation{\affGatech}
    \email{dwayo3@gatech.edu}

    \author{Priscilla Goliat Capriles}
    \affiliation{\fora}

    \author{Masoud Darvish Ganji}
    \affiliation{\ganj}


    \date{\today}
\begin{abstract}
We present SBUP$^3$, a reproducible multiscale workflow that links periodic DFT
dielectric response, reduced multi-transition semiconductor Bloch equations
(SBE), and 2D unidirectional pulse propagation (UPPE) for ultrafast GaAs
simulation. Calibration to the primary Jellison reference selects
$N=1.7585\times10^{28}~\mathrm{m^{-3}}$ and oscillator normalization 0.9593,
with finite 68\% and 95\% uncertainty regions. Post-calibration validation gives
shape RMSE values of 0.3602 for $|\alpha|$ and 0.2872 for $n$ on the primary set
(84 points), and 0.4171 and 0.2944 on an independent secondary Aspnes set
(14 points). Stability checks across $(\Delta t,\Delta z)$ sweep points show
nearly unchanged shape error, and an independent linearity test gives
$\mathrm{ratio}_{P_{\max}}=1.999999430$. UPPE diagnostics show an output spectral
centroid increase (+0.01975 eV) and broadening (+0.00905 eV). These results
establish SBUP$^3$ as a reproducible baseline with explicit provenance, bounded
calibration uncertainty, and stable numerical behavior in the tested regime.
Remaining error is dominated by model-form optical-shape mismatch rather than
solver instability, providing a clear baseline for targeted physics extensions.
\end{abstract}

\maketitle

\section{Introduction}
Ultrafast laser-matter interaction in semiconductors has entered a regime where
sub-cycle field control, coherent interband dynamics, and nonperturbative carrier motion can all be measured directly. Over the last decade, strong-field experiments in solids have demonstrated high-harmonic emission, attosecond-scale band-structure response, and coherent waveform-dependent transport signatures that were previously discussed mainly in atomic and molecular contexts \cite{ghimire2011hhg,schubert2014subcycle,schiffrin2013opticalfield,schultze2013controlling,schultze2014attosecond,luu2015extreme,vampa2015linking,hohenleutner2015realtime,garg2016multipetahertz,ndabashimiye2016carrier,you2017anisotropic,yoshikawa2017graphene,sivis2017tailored,higuchi2017lightfield,langer2018lightwave}. These developments have elevated semiconductors from passive optical media to active ultrafast platforms, with implications for petahertz electronics, compact frequency conversion, and coherent lightwave control
\cite{kruchinin2018strong,ghimire2019highharmonicsolids,goulielmakis2022extreme,heide2024ultrafast}.

This experimental progress has exposed a central modeling challenge: physically
relevant observables are generated across multiple coupled scales. Electronic
coherences and populations evolve at the microscopic level, while measurable
pulse reshaping, spectral transfer, and spatial transport emerge at mesoscopic
or device-relevant scales. If one solves only the microscopic problem, one
obtains detailed polarization physics without realistic propagation context. If
one solves only the propagation problem, one obtains numerically stable field
evolution but often with phenomenological constitutive laws that hide the
underlying electron dynamics. The result is a persistent scale gap in which
parameters are transferred between models in an ad hoc way, usually without
clear uncertainty accounting or reproducible calibration.

Several mature theoretical lines address parts of this problem. Many-body
semiconductor theory and the semiconductor Bloch equations (SBE) provide a
rigorous route from microscopic Hamiltonians to coherent polarization and
population dynamics under ultrafast excitation
\cite{rossi2002ultrafast,axt2004femtosecond}. In the strong-field limit, SBE
and related band-structure formalisms have been used to interpret solid-state
harmonic generation and interband-intraband interference mechanisms
\cite{golde2008hhg,vampa2014theoretical,vampa2015semiclassical,tancogne2017ellipticity}. In parallel,
real-time time-dependent density functional theory (TDDFT) has enabled
first-principles simulation of electron dynamics in crystalline solids under
intense fields, including direct access to transition energies, oscillator
strength trends, and nonlinear response channels
\cite{hohenberg1964inhomogeneous,kohn1965selfconsistent,runge1984densityfunctional,yabana2012tddft,floss2018abinitio}. These approaches are powerful, but
they are not automatically interchangeable: each has distinct numerical
assumptions, response windows, and scaling conventions.

Electromagnetic propagation theory has evolved similarly. Envelope and
unidirectional formulations, including variants of the unidirectional pulse
propagation equation (UPPE), have become standard tools for broadband
ultrashort-pulse transport in dispersive and nonlinear media
\cite{brabec1997nonlinear,kolesik2002nonlinear,kolesik2004unidirectional}.
These methods are computationally efficient and physically transparent at the
field level, but their predictive quality depends critically on how material
polarization is supplied. In semiconductor applications, this coupling step is
often the weak link: microscopic models may output polarization in one set of
units and conventions, while propagation solvers require a macroscopic source
term with consistent scaling, bandwidth support, and numerical regularity.

The literature therefore motivates an integrative viewpoint. Experiments now
probe regimes where microscopic quantum response and propagation dynamics cannot
be decoupled cleanly. At the same time, computational workflows remain fragmented:
first-principles calculations, reduced quantum kinetics, and propagation solvers
are frequently run as separate stages without closed-loop correction. Even when
all three are present, data exchange rules are rarely documented in a way that
supports independent reproduction and calibration. This limits comparability
across studies and slows transition from qualitative interpretation to
quantitative prediction.

%\paragraph{Problem statement.}
The specific problem addressed in this work is the absence of a transparent,
physics-calibrated, and numerically robust workflow for semiconductor ultrafast
dynamics that jointly satisfies three requirements: (i) first-principles
provenance for key material parameters, (ii) tractable microscopic time-domain
carrier and polarization evolution, and (iii) self-consistent propagation-level
field updates under feedback. Existing single-scale studies provide essential
insight but do not by themselves resolve the mapping problem between TDDFT
outputs, DFT-derived transition-reduction outputs, SBE state variables, and
UPPE source terms. In practical terms, this
appears as sensitivity to scaling factors, instability in susceptibility
extraction when the driving spectrum is weak, and insufficient reporting of
convergence behavior across coupling iterations.

%\paragraph{Motivation.}
Our motivation is methodological and translational. Methodologically, we seek a
workflow that preserves physical meaning at every interface: the periodic DFT
dielectric and transition-reduction stage should inform model parameters rather
than merely inspire them; SBE should provide structured nonequilibrium response
rather than black-box fitting; and UPPE should propagate fields with source terms
whose provenance is explicit and testable. Translationally,
this is necessary for publication-quality predictive claims in laser processing,
integrated photonics, and coherent semiconductor control, where reproducibility
depends on whether another group can rerun the same pipeline and recover the same
trends under controlled perturbations. We therefore treat calibration,
robustness, and reporting discipline as central scientific outputs, not
post-processing details.

%\paragraph{Research questions.}
This study is guided by the following research questions:
\begin{enumerate}
\item Can a DFT $\rightarrow$ SBE $\rightarrow$ UPPE pipeline be assembled so
that parameter transfer and unit conventions remain explicit and reproducible?
\item Which numerical controls dominate the reliability of extracted optical
response functions, especially susceptibility-derived absorption in weak-support
spectral regions?
\item Under practical feedback mixing and normalization rules, does the coupled
SBE-UPPE loop converge robustly, and how fast?
\item After calibration to external optical references (primary Jellison, secondary
Aspnes) with a DFT-derived transition basis, what level of
shape agreement and spectral-support coverage is achievable without sacrificing
propagation consistency?
\end{enumerate}

%\paragraph{Objectives.}
To answer these questions, we pursue five concrete objectives:
\begin{enumerate}
\item Build a modular computational workflow with documented file-level data
contracts across DFT, SBE, and UPPE stages.
\item Implement physically interpretable derived outputs ($\alpha$, $n$, and
$\kappa$) from complex susceptibility with explicit support masks.
\item Quantify numerical robustness through targeted sweeps of time-step and
response-support controls.
\item Evaluate coupled-loop stability using reproducible convergence criteria and
feedback diagnostics.
\item Benchmark response-shape agreement against external optical references
using consistent overlay metrics for primary and secondary datasets.
\end{enumerate}

\paragraph{Novel contributions.}
The work makes six contributions relative to common single-scale practice.
First, it introduces SBUP$^3$ as a practical multiscale workflow architecture
rather than a single solver, emphasizing provenance and interface consistency.
Second, it provides an explicit response-support extraction method for
susceptibility inversion that suppresses divide-by-small-field artifacts while
preserving physically relevant bandwidth. Third, it reports calibration and
stability metrics in a machine-readable validation framework, enabling direct
comparison across runs. Fourth, it demonstrates that support-aware tuning can
substantially improve reference-aligned shape fidelity and overlap-point density
without changing global propagation trends. Fifth, it reports coupled-feedback
diagnostics under under-relaxed mixing and explicit normalization, making
scale-alignment limits visible in a reproducible way. Sixth, it positions
sequential and coupled outputs
as complementary diagnostics: the sequential mode as a conservative calibration
baseline and the coupled mode as a higher-gain regime requiring tighter absolute
scaling checks.


\begin{figure*}[t]
\centering
\includegraphics[width=0.8\textwidth]{ study_20260520/F0_intro_gaas_interaction.pdf}
\caption{
Conceptual ultrafast laser pulse drives
nonequilibrium polarization and population dynamics in bulk GaAs, and the
resulting optical response is reported through refractive-index, absorption, and
propagated-field observables.
}
\label{fig:intro_gaas_interaction}
\end{figure*}

Beyond these headline contributions, the conceptual value of this framework lies
in how it reorganizes modeling practice. In much of the ultrafast-solid
literature, simulation components are selected by local convenience: one model
for material properties, another for temporal response, another for propagation,
with manual parameter passing between scripts. This approach can still produce
important qualitative insight, but it makes it difficult to establish whether
agreement with experiment is robust or accidental. A methods platform is useful
only if each interface can be audited. SBUP$^3$ is explicitly designed around
that requirement. Every stage emits reusable arrays with defined units, every
calibration choice is represented as a parameterized control, and every run can
be summarized by machine-readable diagnostics. This is a practical step toward
making multiscale ultrafast modeling falsifiable in the same way as single-stage
numerical benchmarks.

The need for this shift is especially strong for semiconductors because the
relevant physics is neither purely perturbative nor fully ionizing over many
experimentally relevant regimes. Coherent interband polarization, intraband
acceleration, dephasing, and field propagation all contribute on overlapping
timescales. If one insists on purely first-principles propagation across full
device scales, the problem becomes computationally prohibitive for systematic
sensitivity studies. If one moves immediately to effective propagation models
without first-principles anchors, one risks parameter degeneracy, where distinct
microscopic assumptions produce similar macroscopic traces. The DFT-informed
SBE reduction used here is intended as a middle path: retain explicit microscopic
structure where it most directly constrains observables, while enabling enough
runtime efficiency to perform stability and calibration sweeps that are essential
for publication-grade confidence.

GaAs is chosen as the benchmark not because it is the final target material, but
because it offers a stringent and interpretable testbed. Its ultrafast response
has been studied across perturbative and strong-field regimes, and its band
structure enables direct comparison between first-principles transition energies
and reduced-model resonance behavior. In this sense, GaAs functions as a
calibration material: if interface consistency fails here, it is unlikely to
hold in more complex heterostructures or defect-rich systems. Conversely, if a
workflow can maintain numerical stability, traceable scaling, and interpretable
agreement in GaAs, it becomes a credible foundation for more ambitious materials
and geometries.

It is equally important to define scope boundaries. This work does not claim a
fully predictive digital twin of semiconductor strong-field response. The present
SBE stage uses a reduced transition structure, and the calibration comparison is
shape-focused rather than absolute-intensity complete. These choices are
intentional at this stage. The objective is to establish a robust pipeline in
which future physical refinements can be inserted without changing the
reproducibility backbone. In other words, this paper prioritizes methodological
closure before maximal model complexity. That ordering is deliberate: without
numerical and interface robustness, additional physics terms can obscure rather
than improve interpretation.

Figure~\ref{fig:intro_gaas_interaction} gives a visual map of how the driving
ultrafast pulse interacts with bulk GaAs in this study and how that interaction
connects to the optical outputs used for calibration and validation.

By framing the study around reproducible calibration rather than isolated
numerical performance, this paper aims to strengthen the methodological basis for
publishable ultrafast semiconductor simulations. The remainder of the manuscript
details related work, model construction, numerical setup, validation protocol,
and quantitative results for the GaAs benchmark, then discusses current limits
and next steps toward predictive multiscale closure.

\section{Related Work and Positioning}
Prior work relevant to SBUP$^3$ can be grouped into three methodological
streams.

\paragraph{Microscopic semiconductor dynamics (SBE and many-body optics).}
The semiconductor Bloch framework and related many-body formalisms provide a
well-established description of coherent polarization dynamics, carrier
population transfer, and dephasing in ultrafast excitation regimes
\cite{rossi2002ultrafast,axt2004femtosecond}. In strong-field contexts, these
models have been central to interpreting interband-intraband interference and
high-harmonic mechanisms in solids
\cite{golde2008hhg,vampa2014theoretical,vampa2015semiclassical}. Their strength
is direct physical interpretability at the level of quantum kinetics. Their
main limitation for direct comparison to measured pulse observables is that
propagation is often externalized or simplified.

\paragraph{First-principles real-time and linear-response TDDFT.}
TDDFT has enabled predictive access to excitation channels and ultrafast
response trends in solids under strong electromagnetic fields
\cite{yabana2012tddft,floss2018abinitio}. In particular, TDDFT can constrain
effective transition frequencies and oscillator strengths that are otherwise
treated as fit parameters in reduced models. However, fully first-principles
multiscale transport remains computationally expensive for broad sweep-based
robustness studies, and practical interface rules from TDDFT outputs to
reduced propagation models are often under-documented.

\paragraph{Pulse propagation and unidirectional field solvers.}
UPPE-type methods and single-cycle propagation theory are mature and
computationally efficient for broadband ultrashort-pulse evolution in dispersive
media \cite{brabec1997nonlinear,kolesik2002nonlinear,kolesik2004unidirectional}.
These models scale well and are ideal for field-level diagnostics, but their
predictive quality depends critically on the constitutive source term. In
semiconductor applications, this source term is frequently phenomenological or
only weakly connected to first-principles material provenance.

\paragraph{Positioning of SBUP$^3$.}
SBUP$^3$ is positioned at the interface of these streams. It is not introduced
as a replacement for TDDFT, SBE, or UPPE individually; instead, it is a
workflow architecture that enforces explicit contracts between them. The novelty
is operational and methodological:
\begin{itemize}
\item first-principles-informed reduced modeling with explicit scalar handoff
$(\omega_0, E_0, f_0, d_{\mathrm{eff}})$;
\item response-support-aware susceptibility inversion to avoid unstable
low-drive spectral division;
\item machine-readable validation artifacts as first-class outputs;
\item optional closed-loop SBE-UPPE feedback with reproducible convergence
criteria.
\end{itemize}
This positioning reflects the manuscript goal: to establish a robust,
publishable multiscale methods pipeline with auditable assumptions, rather than
to claim final microscopic completeness in a single model stage.

\section{Methodology}
\begin{algorithm}[t]
\caption{SBUP$^3$ end-to-end workflow (release protocol)}
\KwIn{Material system, numerical controls, calibration and coupling flags}
\KwOut{Validated observables and reproducible run artifacts}
Generate or pin periodic dielectric reference and transition arrays\;
Run the SBE stage with release controls and export
$E(t)$, $P(t)$, $n(t)$, $\chi(\omega)$, $\alpha(\omega)$, $n(\omega)$\;
Lift SBE traces to UPPE source terms and propagate
to $z_{\mathrm{end}}$\;
\If{coupling enabled}{
  iterate the SBE-UPPE coupling stage until relative $L^2$ residual
  is below tolerance\;
}
Run calibration/validation and write machine-readable JSON reports\;
Render figure set F0--F6, S1--S4 from pinned reports only\;
\end{algorithm}

\subsection{Periodic DFT Optical Reference and Transition Reduction}
The periodic DFT stage constructs a converged bulk GaAs ground state and then
computes linear optical response on a denser periodic $k$-mesh. This stage
produces an energy-resolved dielectric reference together with derived optical
observables $\{n,\kappa,\alpha\}$. Publication mode enforces non-$\Gamma$-only
sampling and records full provenance metadata for the ground-state and response
steps. In the frozen release used here, the recorded setup is PW/PBE with 32
ground-state bands and 200 eV cutoff, then response evaluation with 48 bands,
$\eta=0.05$ eV, 150 eV response cutoff, and nonlinear frequency mesh up to
8.0 eV.

Reduced transitions are then extracted from the dielectric reference. Peaks are
selected from $\varepsilon_2^{\mathrm{lf}}(\omega)$ with minimum relative height 0.01,
minimum separation 0.02 eV, and bounded cardinality (16--24 transitions). The
release transition model contains 16 transitions in
1.482--2.964 eV with oscillator-sum 0.959264.

\begin{algorithm}[t]
\caption{Periodic DFT dielectric-reference and transition extraction}
\KwIn{Converged periodic ground-state wavefunction artifact}
\KwOut{Dielectric reference, reduced transition set, provenance metadata}
Load periodic ground-state wavefunctions\;
Restart at fixed density on dense response $k$-grid\;
Diagonalize unoccupied bands for linear response\;
Compute dielectric response over nonlinear frequency mesh\;
Convert complex dielectric function to $(n,\kappa,\alpha)$\;
Export optical constants and response metadata\;
Extract reduced transitions from $\varepsilon_2^{\mathrm{lf}}$ peak structure and
export transition arrays with metadata\;
\end{algorithm}

\begin{figure*}[t]
\centering
\includegraphics[width=\textwidth]{ study_20260520/F2_reference_optics_overlay.pdf}
\caption{
Reference optical context in GaAs: dielectric-reference, Jellison, and Aspnes
overlays for absorption and refractive index.
}
\label{fig:method_reference_context}
\end{figure*}

\begin{figure*}[t]
\centering
\includegraphics[width=1\textwidth]{ study_20260520/F3_transition_model.pdf}
\caption{
Reduced transition model extracted from the dielectric reference: selected
transition energies, oscillator weights, and cumulative oscillator sum.
}
\label{fig:method_transition_model}
\end{figure*}

\subsection{Reduced Multi-Transition SBE Dynamics}
The SBE stage evolves transition-resolved coherence
$p_j$ and population $n_j$ for each extracted transition $\omega_j,d_j,f_j$. The
driving field is a Gaussian-envelope carrier
\begin{align}
E_{\mathrm{int}}(t)=E_{\mathrm{amp}}
\exp\!\left[-\left(\frac{t-t_0}{\tau}\right)^2\right]
\cos\!\left[\Omega_L(t-t_0)\right],
\end{align}
optionally mixed with UPPE feedback:
\begin{align}
E(t)=(1-\alpha_{\mathrm{mix}})E_{\mathrm{int}}(t)+
\alpha_{\mathrm{mix}}E_{\mathrm{fb}}(t),
\end{align}
then clipped to $|E|\le E_{\max}$.

The reduced equations are
\begin{align}
\frac{dp_j}{dt} &= -\left(i\omega_j+\gamma_j\right)p_j
+ i\frac{d_j}{\hbar}E(t)\left(1-2n_j\right),\\
\frac{dn_j}{dt} &= -2\,\mathrm{Im}\!\left[\frac{d_j}{\hbar}E(t)p_j^*\right],
\end{align}
integrated with explicit RK4. With normalized weights
$w_j=f_j/\sum_k f_k$, the exported macroscopic polarization is
\begin{align}
P_{\mathrm{mac}}(t)=S_P N \sum_j w_j d_j p_j(t),
\end{align}
where $N=\texttt{SBUP3\_DENSITY\_M3}$ and
$S_P=\texttt{SBUP3\_POLARIZATION\_SCALE}$.

SBUP$^3$ computes both drive-response and dedicated probe-response optics. The
probe branch uses a weak Gaussian pulse
(\texttt{SBUP3\_PROBE\_*} controls) and is used as the primary baseline source in
calibration/validation when available.

Susceptibility inversion is performed on Hann-windowed, mean-subtracted traces:
\begin{align}
\chi(\omega)=\frac{\tilde{P}(\omega)}{\varepsilon_0\tilde{E}(\omega)},
\end{align}
only on trusted support points satisfying
\begin{align}
|\tilde{E}_j| &> \max\!\left(10^{-30},\,\epsilon_{\chi}\max_k|\tilde{E}_k|\right),\\
\sum_{k\le j}\frac{|\tilde{E}_k|^2}{\sum_\ell |\tilde{E}_\ell|^2} &\le p_{\mathrm{cut}},\\
\omega_{\min}\le \omega_j &\le \omega_{\max}.
\end{align}
Optical outputs are then derived from $\tilde{n}=\sqrt{1+\chi}$:
$n=\Re \tilde{n}$, $\kappa=\Im \tilde{n}$, $\alpha=2\omega\kappa/c$.

\begin{algorithm}[t]
\caption{SBE solve with drive/probe optical extraction}
\KwIn{Transition arrays, drive/probe controls, scaling controls}
\KwOut{$E(t)$, $P(t)$, $n(t)$, drive/probe $\chi(\omega)$ and optics outputs}
Construct internal drive waveform\;
\If{feedback available and valid}{
  mix internal and feedback drives with $\alpha_{\mathrm{mix}}$\;
}
Clip field amplitude to $E_{\max}$\;
Initialize $(p_j,n_j)\leftarrow(0,0)$ for all transitions\;
\For{time steps}{
  evaluate reduced multi-transition right-hand side\;
  advance $(p_j,n_j)$ with RK4\;
  accumulate weighted $P_{\mathrm{mac}}(t)$ and $n(t)$\;
}
Extract drive-response $\chi(\omega)$ on trusted support\;
\If{probe mode enabled}{
  run weak probe solve and extract probe-response $\chi(\omega)$\;
}
Write drive/probe optics arrays and transition metadata\;
\end{algorithm}

\subsection{UPPE Propagation Model}
The UPPE stage is solved on a Cartesian $(x,t)$ grid with propagation along
$z$. SBE temporal traces are lifted to 2D by a Gaussian transverse envelope,
producing initial field $E(x,t,0)$ and source polarization $P(x,t)$. The solver
uses a spectral representation in $(k_x,\omega)$ with a unidirectional update
operator using $k_z=\sqrt{\max(k_0^2-k_x^2,0)}$, where
$k_0(\omega)=n(|\omega|)\omega/c$ from a GaAs dispersion model.

The propagation step combines exact linear phase accumulation
$\exp(-ik_z\Delta z)$ with an additive source term proportional to
$\omega^2\hat{P}/(2k_z+\varepsilon)$. The implementation is intentionally simple
and stable for the present benchmark regime; polarization is treated as fixed
within each UPPE run and updated only through outer-loop coupling.

\begin{algorithm}[t]
\caption{UPPE propagation in $(k_x,\omega)$ domain}
\KwIn{$E_t$, $P_t$, $x$-grid, $t$-grid, $w_0$, $\Delta z$, $N_z$}
\KwOut{$E(x,t,z_{\mathrm{end}})$}
Lift 1D traces to 2D with Gaussian transverse envelope\;
FFT to obtain $\hat{E}(k_x,\omega)$ and $\hat{P}(k_x,\omega)$\;
Build dispersion arrays $k_0(\omega)$ and $k_z(k_x,\omega)$\;
\For{$m=1$ to $N_z$}{
  apply linear propagator $e^{-ik_z\Delta z}$\;
  add source update
  $i\Delta z\,\omega^2\hat{P}/(2k_z+\varepsilon)$\;
}
Inverse FFT and take real part\;
\end{algorithm}

\subsection{Self-Consistent Coupling Loop}
In coupled mode, SBUP$^3$ iterates between SBE and UPPE until feedback
stabilizes. After each UPPE run, the on-axis field
$E_{\mathrm{out}}(t)=E(x\approx0,t)$ is extracted and combined with the previous
feedback waveform through under-relaxed mixing. Optional RMS or peak
normalization rescales feedback to the current SBE drive amplitude before
writing the feedback waveform for the next iteration.

Convergence is monitored by a relative $L^2$ update norm with absolute floor:
\begin{align}
r_q=\frac{\|E_{\mathrm{fb}}^{(q)}-E_{\mathrm{fb}}^{(q-1)}\|_2}
{\max(\|E_{\mathrm{fb}}^{(q-1)}\|_2,\epsilon_{\mathrm{abs}})}.
\end{align}
Iteration stops once $r_q<r_{\mathrm{tol}}$ or once the iteration cap is
reached.

\begin{algorithm}[t]
\caption{Self-consistent SBE-UPPE coupling}
\KwIn{Initial controls, $\beta$, normalization mode, $r_{\mathrm{tol}}$, $N_{\max}$}
\KwOut{Converged feedback field and final propagated state}
Set previous feedback to null\;
\For{$q=1$ to $N_{\max}$}{
  run SBE with current feedback file (or internal drive at first pass)\;
  run UPPE using current SBE outputs\;
  extract on-axis output waveform $E_{\mathrm{out}}^{(q)}(t)$\;
  \uIf{$q=1$}{
    set $E_{\mathrm{fb}}^{(1)}\leftarrow E_{\mathrm{out}}^{(1)}$\;
  }\Else{
    set $E_{\mathrm{fb}}^{(q)}\leftarrow
    (1-\beta)E_{\mathrm{fb}}^{(q-1)}+\beta E_{\mathrm{out}}^{(q)}$\;
  }
  optionally normalize feedback in RMS or peak mode\;
  write feedback file for next SBE pass\;
  \If{$q>1$ and relative $L^2$ update $< r_{\mathrm{tol}}$}{
    stop iteration\;
  }
}
\end{algorithm}

\subsection{Physical Calibration and Uncertainty Model}
Calibration is performed in physical mode, which fits carrier density $N$ and oscillator normalization
$s_{\mathrm{osc}}$ (rather than a free global scale only). Let channel-wise
comparison metrics be computed over overlapping energy windows using normalized
RMSE. The data objective is
\begin{align}
\mathcal{J}_{\mathrm{data}}=
\frac{\sum_c w_c\,\mathrm{RMSE}^{(c)}_{\mathrm{norm}}}{\sum_c w_c},
\end{align}
for channels $c\in\{\alpha,n,\chi',\chi''\}$ selected by the user. A log-space
Gaussian prior penalty is added:
\begin{align}
\mathcal{J}_{\mathrm{tot}}=
\mathcal{J}_{\mathrm{data}}+\lambda_{\mathrm{prior}}
\left[\Pi_N(N)+\Pi_s(s_{\mathrm{osc}})\right].
\end{align}
The release calibration used equal channel weights for $\alpha$ and $n$,
$\lambda_{\mathrm{prior}}=0.05$, and 101$\times$101 log grids over
$N\in[0.1,10]\times N_{\mathrm{prior}}$ and
$s_{\mathrm{osc}}\in[0.5,2.0]$.

Uncertainty is reported as objective-threshold surrogate regions:
68\% and 95\% boxes are defined by
$\mathcal{J}_{\mathrm{tot}}\le \mathcal{J}_{\min}+\Delta$ with
$(\Delta_{68},\Delta_{95})=(0.02,0.05)$. This is used as the manuscript
uncertainty summary for calibration identifiability and residual shape mismatch.

\section{Numerical Setup and Reproducibility}
\subsection{Software Environment}
The release workflow uses documented software environments for each stage. The
primary run/calibration/validation environment records Python 3.14.0 and NumPy
2.4.5. The periodic DFT stage records GPAW 25.7.0 with Python 3.10.18, ASE
3.28.0, NumPy 2.2.6, and SciPy 1.15.3. Figure rendering is performed with
Matplotlib 3.10.9. The environment snapshot and pinned export values are stored
in the primary release-freeze record. Execution, validation, calibration, and
figure generation are performed through dedicated workflow scripts that read
only pinned study report paths.

\subsection{Baseline Configuration}
For the frozen release baseline (sequential mode), controls from
the primary release run record are:
$T=50$ fs, $\Delta t=5$ as, $\gamma=5\times10^{13}~\mathrm{s^{-1}}$,
$E_{\mathrm{amp}}=10^7$ V/m, $\Omega_L=9.9883\times10^{14}$ rad/s,
$E_{\max}=5\times10^7$ V/m, density
$N=1.7585466510231837\times10^{28}~\mathrm{m^{-3}}$,
$s_{\mathrm{osc}}=0.9592641193252643$, and
$S_P=\texttt{SBUP3\_POLARIZATION\_SCALE}=1$.

Trusted support controls are
$\epsilon_{\chi}=10^{-12}$, $p_{\mathrm{cut}}=1.0$,
$E_{\max}^{(\chi)}=3.0$ eV, with probe window
1.5--3.0 eV and probe amplitude $5\times10^4$ V/m.
UPPE controls are $N_x=256$, $x$-window $10~\mu$m, $w_0=2~\mu$m,
$\Delta z=1~\mu$m, $N_z=100$.
Coupling defaults (used in separate stability checks) are
$\beta=0.2$, RMS normalization, $r_{\mathrm{tol}}=10^{-3}$, and
$N_{\max}=10$.

\begin{table}[t]
\centering
\caption{Primary simulation controls used for reported runs.}
\label{tab:numerical_controls}
\begin{tabular}{ll}
\toprule
Parameter & Value \\
\midrule
$T$ (SBE window) & $50~\mathrm{fs}$ \\
$\Delta t$ (SBE) & $5~\mathrm{as}$ \\
$\gamma$ & $5\times10^{13}~\mathrm{s^{-1}}$ \\
$E_{\mathrm{amp}}$ & $10^7~\mathrm{V/m}$ \\
$E_{\max}$ (clip) & $5\times10^7~\mathrm{V/m}$ \\
$N$ (density) & $1.758546651\times10^{28}~\mathrm{m^{-3}}$ \\
$s_{\mathrm{osc}}$ & $0.9592641193$ \\
$S_P$ & $1.0$ \\
$E_{\mathrm{probe}}$ & $5\times10^4~\mathrm{V/m}$ \\
$[E_{\min}^{\chi,\mathrm{probe}},E_{\max}^{\chi,\mathrm{probe}}]$ & $[1.5,3.0]~\mathrm{eV}$ \\
$N_x$ (UPPE) & $256$ \\
$x$ window & $10~\mu\mathrm{m}$ \\
$w_0$ & $2~\mu\mathrm{m}$ \\
$\Delta z$, $N_z$ & $1~\mu\mathrm{m}$, $100$ \\
$\epsilon_{\chi}$ & $10^{-12}$ \\
$p_{\mathrm{cut}}$ & $1.0$ \\
$E_{\max}^{(\chi)}$ & $3.0~\mathrm{eV}$ \\
\bottomrule
\end{tabular}
\end{table}

\subsection{Validation Protocol}
Validation is executed by the automated validation stage. Each report
stores finite-value counts, spectral centroids/RMS widths, dephasing-tail
statistics, transverse energy-proxy drift, support coverage fractions, baseline
fit metrics (scale, RMSE, normalized RMSE, shape RMSE, overlap count), and
optional linearity checks.

Primary baseline validation uses Jellison (1992) data in the probe-response
window 1.5716--2.9777 eV. Secondary validation uses Aspnes (1986) on the same
window. The release linearity check reports
$\mathrm{ratio}_{P_{\max}}=1.999999430$ at
$E_{\mathrm{amp}}=10^7$ V/m (ideal value 2.0).

The release convergence panel (S4A) is built from four explicit sweep points:
$(\Delta t,\Delta z)\in\{(5~\mathrm{as},1.0~\mu\mathrm{m}),
(5~\mathrm{as},0.5~\mu\mathrm{m}),
(2.5~\mathrm{as},1.0~\mu\mathrm{m}),
(2.5~\mathrm{as},0.5~\mu\mathrm{m})\}$,
with reported $\alpha_{\mathrm{abs}}$ shape RMSE, $n$ shape RMSE, and support
coverage per point.

\paragraph{Uncertainty statement for residual shape mismatch.}
The remaining mismatch is quantified rather than inferred qualitatively:
for the primary release baseline,
$\mathrm{RMSE}^{\mathrm{shape}}_{\alpha_{\mathrm{abs}}}=0.3602$ and
$\mathrm{RMSE}^{\mathrm{shape}}_{n}=0.2872$ over 84 overlap points, while the
physical calibration uncertainty ranges are
$N\in[1.6038,2.0191]\times10^{28}~\mathrm{m^{-3}}$ and
$s_{\mathrm{osc}}\in[0.8467,1.1019]$ (68\%) and
$N\in[1.4627,2.2139]\times10^{28}~\mathrm{m^{-3}}$,
$s_{\mathrm{osc}}\in[0.7792,1.1975]$ (95\%). These bounds indicate that the
dominant residual error is model-form shape mismatch, not numerical instability.

\begin{algorithm}[t]
\caption{Validation and tuning protocol}
\KwIn{Reference dielectric-response dataset, baseline controls, sweep grids}
\KwOut{Selected operating point and validation reports}
Run baseline sequential pipeline and compute validation metrics\;
Compute dielectric-reference overlay metrics (scale, RMSE, shape RMSE, overlap count)\;
Execute coarse sweep over temporal and response-support controls\;
Rank candidates by shape RMSE and overlap-point density\;
Execute focused sweep around top candidates\;
Select operating point with best shape-support tradeoff\;
Run coupled verification and check feedback convergence\;
Archive reports from selected controls\;
\end{algorithm}

\section{Results}
This section reports four result layers in sequence: calibration-point
identification, transfer to an independent reference set, numerical robustness,
and propagation diagnostics. This ordering separates parameter fitting from
dataset shift and discretization effects.

\subsection{Release Calibration Point}
The frozen physical-fit calibration (Jellison primary reference) selects
$N=1.758546651\times10^{28}~\mathrm{m^{-3}}$ and
$s_{\mathrm{osc}}=0.9592641193$, with effective ratio
$r_{\mathrm{best}}=0.9549925860$. The objective decomposition is
$\mathcal{J}_{\mathrm{tot}}=0.55484$,
$\mathcal{J}_{\mathrm{data}}=0.55097$,
$\mathcal{J}_{\mathrm{prior}}=0.07733$ (prior weight $0.05$).
The calibrated channel metrics over 1.5716--2.9777 eV are
$\mathrm{RMSE}_{\mathrm{norm}}(\alpha)=0.4817$ and
$\mathrm{RMSE}_{\mathrm{norm}}(n)=0.6202$ (84 overlap points).

Figure~\ref{fig:results_calibration} shows the objective landscape and the
surrogate uncertainty envelopes. The reported 68\% set is
$N\in[1.6038,2.0191]\times10^{28}~\mathrm{m^{-3}}$,
$s_{\mathrm{osc}}\in[0.8467,1.1019]$; the 95\% set is
$N\in[1.4627,2.2139]\times10^{28}~\mathrm{m^{-3}}$,
$s_{\mathrm{osc}}\in[0.7792,1.1975]$.

The uncertainty sets indicate a bounded neighborhood rather than a flat
degenerate plateau. The data term dominates the total objective, while the prior
term regularizes extreme density-oscillator combinations. We therefore treat
this calibration point as the operational baseline for transfer and robustness
tests, not as a claim of complete microscopic closure.

\begin{figure*}[t]
\centering
\includegraphics[width=\textwidth]{ study_20260520/F4_calibration_landscape.pdf}
\caption{
Release physical-fit calibration landscape and surrogate uncertainty bounds from
the pinned calibration report.
}
\label{fig:results_calibration}
\end{figure*}

\subsection{Post-Calibration Optical Agreement}
Primary validation against Jellison (probe-response baseline) gives
$\mathrm{RMSE}_{\mathrm{norm}}(\alpha_{\mathrm{abs}})=0.2470$,
$\mathrm{RMSE}^{\mathrm{shape}}(\alpha_{\mathrm{abs}})=0.3602$,
$\mathrm{RMSE}_{\mathrm{norm}}(n)=0.2600$, and
$\mathrm{RMSE}^{\mathrm{shape}}(n)=0.2872$ over 84 overlap points in
1.5716--2.9777 eV. The fitted amplitude factor from this baseline comparison is
1.1718 for $\alpha_{\mathrm{abs}}$.

Secondary validation against Aspnes gives
$\mathrm{RMSE}_{\mathrm{norm}}(\alpha_{\mathrm{abs}})=0.2575$,
$\mathrm{RMSE}^{\mathrm{shape}}(\alpha_{\mathrm{abs}})=0.4171$,
$\mathrm{RMSE}_{\mathrm{norm}}(n)=0.2629$, and
$\mathrm{RMSE}^{\mathrm{shape}}(n)=0.2944$ over 14 overlap points, with fitted
factor 1.0408. The lower overlap cardinality in the secondary set is consistent
with its sparser sampled reference grid in the release window.

The transfer result is consistent with moderate degradation under dataset shift.
Refractive-index metrics remain close between primary and secondary references,
while absorption-shape mismatch increases on the secondary set. The secondary
overlap count is smaller (14 points), so local reference differences carry more
weight in aggregate metrics. The fitted absorption scale changes from 1.1718
(primary) to 1.0408 (secondary), indicating bounded redistribution between scale
and shape rather than loss of overall spectral structure.

\begin{figure*}[t]
\centering
\includegraphics[width=0.7\textwidth]{  study_20260520/F5_postcal_fit_jellison.pdf}
\caption{
Post-calibration primary overlay (Jellison): simulated probe-response
$|\alpha|$ and $n$ against reference data in the 1.5--3.0 eV fit window.
}
\label{fig:results_postcal}
\end{figure*}

\subsection{Stability and Linearity}
The convergence check uses four explicit sweep points:
$(\Delta t,\Delta z)=(5~\mathrm{as},1.0~\mu\mathrm{m})$,
$(5~\mathrm{as},0.5~\mu\mathrm{m})$,
$(2.5~\mathrm{as},1.0~\mu\mathrm{m})$,
$(2.5~\mathrm{as},0.5~\mu\mathrm{m})$.
Across these points,
$\mathrm{RMSE}^{\mathrm{shape}}(\alpha_{\mathrm{abs}})$ remains in
[0.36019, 0.36021], while
$\mathrm{RMSE}^{\mathrm{shape}}(n)$ remains in
[0.28718, 0.28791]. The $\Delta z$ change has no visible effect at reported
precision; changing $\Delta t$ mainly changes support coverage
(probe support 0.003599 at 5 as vs 0.001800 at 2.5 as) with nearly unchanged
shape error.

Linearity is verified at $E_{\mathrm{amp}}=10^7$ V/m with
$\mathrm{ratio}_{P_{\max}}=1.999999430$, i.e., relative deviation
$|\mathrm{ratio}-2|/2=2.85\times10^{-7}$ from ideal first-order behavior.
Figure~\ref{fig:results_stability} summarizes both checks.

Within the tested control rectangle, the optical-shape conclusions are unchanged.
Changing $\Delta z$ from 1.0 to 0.5~$\mu$m does not materially affect reported
shape metrics. Changing $\Delta t$ from 5 to 2.5~as mainly reduces trusted
support coverage, with nearly unchanged shape error.

The linearity check supports this robustness result. The reported ratio is
effectively ideal at the chosen amplitude, showing that response extraction and
polarization scaling are consistent with local first-order behavior in this
baseline regime.

\begin{figure*}[t]
\centering
\includegraphics[width=\textwidth]{ study_20260520/S4_stability_linearity.pdf}
\caption{
Release numerical diagnostics: explicit convergence values across
$(\Delta t,\Delta z)$ combinations and independent linearity check.
}
\label{fig:results_stability}
\end{figure*}

\subsection{Propagation Diagnostics}
For the release baseline, spectral moments from validation give
$(\mu_{\mathrm{in}},\sigma_{\mathrm{in}})=(0.65744,0.09746)$ eV and
$(\mu_{\mathrm{out}},\sigma_{\mathrm{out}})=(0.67720,0.10651)$ eV, corresponding
to a $+0.01975$ eV centroid shift and $+0.00905$ eV width increase at the UPPE
output. Thus, in this release configuration, propagation broadens the spectrum
and slightly shifts its centroid upward.

To monitor propagation health, we track the transverse energy-like proxy
\begin{align}
Q_{\mathrm{proxy}}(t)=\int |E(x,t)|^2\,dx.
\end{align}
This quantity remains finite throughout the run. Its relative drift and
standard-deviation indicators are $\mathrm{drift}_{\mathrm{rel}}=6.7456$ and
$\mathrm{std}_{\mathrm{rel}}=1.4600$.

The positive centroid shift and broader output width indicate measurable spectral
redistribution during propagation. The shift is small but systematic, consistent
with controlled nonlinear transport in the tested configuration.

The $Q_{\mathrm{proxy}}$ indicators should be read as numerical health monitors,
not strict conservation invariants. In the present reduced closure, the material
source term is injected through a simplified source-updated path, so exact
energy conservation is not the target diagnostic in this section. The key point
is that the fields remain finite, the proxy remains trackable, and the
spectral-moment trends are stable enough for cross-run comparison and future
model refinement.

Taken together, the Results show a bounded calibration region, transfer with
moderate degradation on an independent reference set, numerical stability across
the tested $(\Delta t,\Delta z)$ points, and interpretable propagation trends.

\begin{figure*}[t]
\centering
\includegraphics[width=\textwidth]{ study_20260520/F6_uppe_field_evolution.pdf}
\caption{
Release UPPE propagation output: spatiotemporal field magnitude and transverse
energy proxy from pinned baseline artifacts.
}
\label{fig:results_uppe_release}
\end{figure*}

\begin{table*}[t]
\centering
\caption{
Pinned release optical-fit summary for primary (Jellison) and secondary
(Aspnes) validations.
}
\label{tab:results_summary}
\begin{tabular}{lcccccc}
\toprule
Reference set & $N_{\mathrm{pts}}$ & $\mathrm{RMSE}_{\mathrm{norm}}(\alpha_{\mathrm{abs}})$ & $\mathrm{RMSE}_{\mathrm{shape}}(\alpha_{\mathrm{abs}})$ & $\mathrm{RMSE}_{\mathrm{norm}}(n)$ & $\mathrm{RMSE}_{\mathrm{shape}}(n)$ & fitted $\alpha_{\mathrm{abs}}$ scale \\
\midrule
Primary (Jellison 1992) & 84 & 0.2470 & 0.3602 & 0.2600 & 0.2872 & 1.1718 \\
Secondary (Aspnes 1986) & 14 & 0.2575 & 0.4171 & 0.2629 & 0.2944 & 1.0408 \\
\bottomrule
\end{tabular}
\setlength{\tabcolsep}{6pt}
\end{table*}

\section{Discussion}
The release study supports three method-level conclusions. First, the workflow
is reproducible: run, calibration, and validation artifacts are pinned and
hash-tracked, and the manuscript figures are regenerated from those exact
artifacts. Second, the selected operating point is numerically stable under the
reported sweep controls: shape metrics remain nearly invariant across tested
$(\Delta t,\Delta z)$ values, and the linearity check remains effectively ideal
($\mathrm{ratio}_{P_{\max}}=1.999999430$). Third, calibration is identifiable:
the best-fit point lies within finite 68\%/95\% objective neighborhoods rather
than on a broad flat direction.

The main residual limitation remains model-form mismatch in optical shape, not
solver instability. This is consistent with the reduced transition basis and the
current single-scale macroscopic mapping. In practical terms, the current model
captures reference-aligned trends in the fitted window, but still leaves
structured residuals that should be attributed to missing physics (broader
multi-band coupling, excitonic structure, and refined damping channels) before
claiming predictive closure outside that calibrated regime.

\begin{table*}[t]
\centering
\caption{
External optical reference datasets used in the pinned release workflow.
}
\label{tab:discussion_reference_sources}
\small
\setlength{\tabcolsep}{4pt}
\begin{tabular}{p{0.23\textwidth}p{0.10\textwidth}p{0.12\textwidth}p{0.18\textwidth}p{0.30\textwidth}}
\toprule
Reference set & Citation & Native range (eV) & Provenance layer & Role in this study \\
\midrule
Jellison GaAs optical functions (1992) &
\cite{jellison1992optical} &
1.476--5.298 &
Primary local reference snapshot &
Primary physical-fit calibration target and primary post-calibration validation \\
Aspnes GaAs/AlGaAs optical properties (1986) &
\cite{aspnes1986optical} &
1.500--6.001 &
Secondary local reference snapshot &
Secondary validation for out-of-sample consistency at fixed calibrated controls \\
RefractiveIndex.INFO source manifests &
\cite{polyanskiy2024refractiveindexinfo} &
n/a &
Source-manifest import layer &
Import provenance layer for local experimental reference datasets \\
\bottomrule
\end{tabular}
\setlength{\tabcolsep}{6pt}
\end{table*}

Table~\ref{tab:discussion_reference_sources} makes the reference-data boundary
explicit: Jellison data define the fitted operating point, Aspnes data test
cross-dataset behavior, and refractiveindex.info provides the import-provenance
container for the local reference snapshots. This separation is important for
interpretation because improved agreement with the primary fit set does not by
itself imply robustness to dataset shift; the secondary reference is therefore a
required check, not an optional visualization.

The next step is to retain this release protocol while adding stronger
absolute-scale checks and richer reduced-transition physics. This preserves
comparability across revisions and improves physical completeness.

\section{Practical Implications}
This framework gives teams a reliable way to test ultrafast semiconductor
scenarios before expensive lab campaigns. Instead of broad trial-and-error
experiments, researchers can narrow parameter ranges in simulation and then
prioritize the most informative measurements.

It also improves day-to-day collaboration. Because the same pinned data, scripts,
and figures can be rerun by different users, students and collaborators can
reproduce the same baseline results and continue from the same starting point.

For the broader community, the main value is a reusable reproducible template:
shared inputs, fixed analysis steps, and transparent reporting of fit quality
and remaining mismatch. This makes cross-study comparison easier and reduces
ambiguity when extending the workflow to related materials.

\section{Conclusion}
SBUP$^3$ now provides a release-frozen, reproducible multiscale workflow that
connects DFT-derived inputs, reduced semiconductor Bloch dynamics, and UPPE
propagation for ultrafast GaAs studies. All reported figures and metrics are
generated from pinned run, calibration, and validation artifacts.

The physical-fit release calibration selects
$N=1.758546651\times10^{28}~\mathrm{m^{-3}}$ and
$s_{\mathrm{osc}}=0.9592641193$
($r_{\mathrm{best}}=0.9549925860$), with finite 68\%/95\% objective regions.
Post-calibration agreement is consistent across references: on the primary
Jellison set, shape RMSE is 0.3602 for $\alpha_{\mathrm{abs}}$ and 0.2872 for
$n$ (84 points), while on the secondary Aspnes set it is 0.4171 and 0.2944
(14 points). Numerical behavior is stable under the reported
$(\Delta t,\Delta z)$ sweeps, and the independent linearity test remains nearly
ideal
($\mathrm{ratio}_{P_{\max}}=1.999999430$). UPPE diagnostics show a small
centroid up-shift ($+0.01975$ eV) with spectral broadening ($+0.00905$ eV).

These results support SBUP$^3$ as a credible methods baseline for reproducible
ultrafast semiconductor modeling. The main remaining limitation is optical
shape mismatch from model form, not solver instability. The most direct next
steps are to strengthen absolute-scale constraints, expand the reduced
transition basis toward multi-band/excitonic response, and keep the same pinned
release discipline so future improvements remain directly comparable.


\section*{Nomenclature}
\begin{table}[!h]
\centering
\caption{Core symbols used in SBUP$^3$.}
\label{tab:nomenclature}
\small
\setlength{\tabcolsep}{2pt}
\begin{tabular*}{\columnwidth}{@{\extracolsep{\fill}}lll}
\toprule
Symbol & Definition & Units \\
\midrule
$E_0$ & Bright excitation energy & eV \\
$\omega_0$ & Transition frequency & rad s$^{-1}$ \\
$f_0$ & Oscillator strength & 1 \\
$d_{\mathrm{eff}}$ & Effective dipole & C m \\
$E(t)$ & SBE drive field & V m$^{-1}$ \\
$E_{\mathrm{fb}}(t)$ & UPPE feedback field & V m$^{-1}$ \\
$p(t)$ & Interband coherence & 1 \\
$n(t)$ & Excited population & 1 \\
$\gamma$ & Dephasing rate & s$^{-1}$ \\
$P_{\mathrm{mac}}(t)$ & Macroscopic polarization & C m$^{-2}$ \\
$\chi(\omega)$ & Susceptibility & 1 \\
$n(\omega),\kappa(\omega)$ & Refractive and loss terms & 1 \\
$\alpha(\omega)$ & Absorption coefficient & m$^{-1}$ \\
$\Delta t,\Delta z$ & Time and z steps & s, m \\
$N_x,N_z$ & Grid points ($x,z$) & 1 \\
$\beta$ & Mixing factor & 1 \\
$r_q,r_{\mathrm{tol}}$ & Residual, tolerance & 1 \\
\bottomrule
\end{tabular*}
\end{table}


\section*{Author Contributions}

D.D.K.W: conceptualization, methodology, validation and visualization, software, writing – original draft, review \& editing. P.G.C: methodology, validation and visualization, software, writing – original draft, review \& editing.  M.D.G: methodology, validation and visualization, software, writing – original draft, review \& editing.

\section*{Acknowledgment(s)}
The authors acknowledge contributors, users, and reviewers who provided feedback on decoding workflows, reproducibility scripts, and documentation quality. Any opinions, findings, conclusions, or recommendations expressed in this research are those of the author(s) and do not necessarily reflect the views of their respective affiliations.


\section*{Data and Code Availability}
Code and analysis workflows are available at: \url{https://github.com/DennisWayo/SBUP3-ultrafast}

\section*{Funding}
This research was not funded.


\section*{Disclosure statement}
No potential conflict of interest was reported by the author(s).


\appendix

\section{Equation Details}\label{sec:appendix_math}
\subsection{Legacy Scalar Mapping (Compatibility Path)}
For compatibility with legacy single-transition runs, the electronic-structure
stage can provide a bright excitation energy $E_0$ (in eV) and oscillator
strength $f_0$. These are mapped to effective two-level SBE parameters as
\begin{align}
\omega_0 &= \frac{E_0 e}{\hbar}, \\
|d_{\mathrm{eff}}| &= \sqrt{\frac{3\hbar e^2}{2m_e\omega_0} f_0},
\end{align}
where $e$ is the elementary charge and $m_e$ is the electron mass. The second
relation is the isotropic oscillator-strength inversion used in the
compatibility code path.

\subsection{Driven Reduced Semiconductor Bloch Model}
The internal driving field is defined as a Gaussian-envelope carrier,
\begin{align}
E_{\mathrm{int}}(t)
= E_{\mathrm{amp}}
\exp\!\left[-\left(\frac{t-t_0}{\tau}\right)^2\right]
\cos\!\left[\Omega_L(t-t_0)\right],
\end{align}
with $t_0=T/2$ and $\tau=0.15T$ for total window $T$.



\begin{table*}[t]
\centering
\caption{
Supplementary release snapshots from pinned machine-readable validation reports
(support fraction corresponds to probe-response support coverage).
}
\label{tab:appendix_validation}
\small
\setlength{\tabcolsep}{3pt}
\begin{tabular}{lccccccc}
\toprule
Case & Report ID & RMSE$_{\mathrm{shape}}(|\alpha|)$ & RMSE$_{\mathrm{shape}}(n)$ & $N_{\mathrm{pts}}$ & Support frac. & Feedback rel-$L^2$ & Drift rel. \\
\midrule
Primary release (Jellison) & 20260520\_081916\_104478 & 0.36021 & 0.28718 & 84 & 0.003599 & 1.425 & 6.746 \\
Convergence ($\Delta t=5$ as, $\Delta z=1.0~\mu$m) & 20260520\_081928\_899769 & 0.36021 & 0.28718 & 84 & 0.003599 & 2.250 & 6.746 \\
Convergence ($\Delta t=2.5$ as, $\Delta z=0.5~\mu$m) & 20260520\_081928\_899769 & 0.36019 & 0.28791 & 84 & 0.001800 & 1.447 & 6.746 \\
Secondary release (Aspnes) & 20260520\_082112\_511494 & 0.41706 & 0.29439 & 14 & 0.001800 & 1.447 & 6.746 \\
\bottomrule
\end{tabular}
\end{table*}
\setlength{\tabcolsep}{6pt}


\begin{table*}[t]
\centering
\caption{Periodic DFT dielectric-response setup used in this work (GaAs benchmark).}
\label{tab:appendix_dft_settings}
\begin{tabular}{ll}
\toprule
Setting & Value \\
\midrule
Code & GPAW + ASE \\
Ground-state XC & PBE \\
$\,$Ground-state mode & PW (plane-wave) \\
$k$-point grid (ground state) & $(4,4,4)$ (frozen metadata) \\
$k$-point grid (response) & $(6,6,6)$ (frozen metadata) \\
PW cutoff (ground state) & 200 eV \\
Number of bands (ground state) & 32 \\
Number of bands (response) & 48 (frozen metadata) \\
Smearing (Fermi-Dirac) & 0.01 eV \\
Dielectric broadening $\eta$ & 0.05 eV \\
Dielectric cutoff & 150 eV \\
Frequency mesh & nonlinear ($d\omega_0=0.02$, $\omega_2=5.0$, $\omega_{\max}=8.0$ eV) \\
Integration mode & point integration \\
Intraband term & disabled (semiconductor response path) \\
Ground-state logs & \texttt{gaas\_pw.txt} \\
Ground-state metadata & \texttt{gaas\_ground\_state\_meta.json} \\
Dielectric metadata & \texttt{gaas\_dielectric\_meta.json} \\
\bottomrule
\end{tabular}
\end{table*}



\begin{table*}[t]
\centering
\caption{Legacy scalar transition outputs (compatibility path, not primary calibration).}
\label{tab:appendix_tddft_outputs}
\begin{tabular}{ll}
\toprule
Quantity & Value \\
\midrule
$E_0$ & 0.657444 eV \\
$\omega_0$ & $9.98833\times 10^{14}$ rad/s \\
$f_0$ & $1.25957\times 10^{-3}$ \\
$|d_{\mathrm{eff}}|$ & $2.37090\times 10^{-30}$ C$\cdot$m \\
$|d_{\mathrm{eff}}|$ (Debye) & 0.711 D \\
Reference transition entries (CSV) & 48 \\
Release primary fit reference (outside compatibility path) & \texttt{analysis/data/gaas\_exp\_jellison\_1992.csv} \\
\bottomrule
\end{tabular}
\end{table*}


If feedback exists and is enabled, the drive is mixed as
\begin{align}
E(t) = (1-\alpha_{\mathrm{mix}})E_{\mathrm{int}}(t)
+ \alpha_{\mathrm{mix}}E_{\mathrm{fb}}(t),
\end{align}
then clipped to $|E(t)|\le E_{\max}$ for stability.

For transition-resolved states $(p_j,n_j)$, the implemented reduced equations
are
\begin{align}
\frac{dp_j}{dt}
&= -\left(i\omega_j+\gamma_j\right)p_j
+ i\frac{d_j}{\hbar}E(t)\left(1-2n_j\right), \\
\frac{dn_j}{dt}
&= -2\,\mathrm{Im}\!\left[
\frac{d_j}{\hbar}E(t)p_j^*(t)\right],
\end{align}
where $\gamma_j$ may include the optional linear energy-dependent slope used in
the code. With normalized oscillator weights
$w_j=f_j/\sum_k f_k$, the effective coherence is
\begin{align}
p_{\mathrm{eff}}(t)=\sum_j w_j d_j p_j(t),
\end{align}
and the macroscopic polarization exported to UPPE is
\begin{align}
P_{\mathrm{mac}}(t)
= S_P N\,p_{\mathrm{eff}}(t),
\end{align}
with density $N$ and user scale factor $S_P$.

\subsection{RK4 Time Integration}
With timestep $\Delta t$, the solver updates the stacked transition state
$\mathbf{y}=(\mathbf{p},\mathbf{n})^T$ using explicit RK4:
\begin{align}
\mathbf{y}_{m+1}
= \mathbf{y}_m + \frac{\Delta t}{6}
\left(\mathbf{k}_1+2\mathbf{k}_2+2\mathbf{k}_3+\mathbf{k}_4\right),
\end{align}
where
$\mathbf{k}_1=\mathbf{F}(\mathbf{y}_m,t_m)$,
$\mathbf{k}_2=\mathbf{F}(\mathbf{y}_m+\frac{\Delta t}{2}\mathbf{k}_1,t_m+\frac{\Delta t}{2})$,
$\mathbf{k}_3=\mathbf{F}(\mathbf{y}_m+\frac{\Delta t}{2}\mathbf{k}_2,t_m+\frac{\Delta t}{2})$,
$\mathbf{k}_4=\mathbf{F}(\mathbf{y}_m+\Delta t\,\mathbf{k}_3,t_m+\Delta t)$.

\subsection{Susceptibility Extraction with Response Support}
After propagation, the code applies a Hann window to mean-subtracted traces:
\begin{align}
\tilde{E}_j &= \mathcal{F}\!\left[w_m(E_m-\bar{E})\right]_j, \\
\tilde{P}_j &= \mathcal{F}\!\left[w_m\mathrm{Re}(P_m-\bar{P})\right]_j.
\end{align}
The frequency grid is the nonnegative branch from a real-input FFT
($\mathrm{rFFT}$), $\omega_j=2\pi f_j$.

Susceptibility is inverted only on a trusted support mask $\mathcal{S}$:
\begin{align}
\chi_j = \frac{\tilde{P}_j}{\varepsilon_0\tilde{E}_j},
\qquad j\in\mathcal{S},
\end{align}
and set to NaN otherwise. The mask is the intersection of four criteria:
\begin{align}
|\tilde{E}_j| &> \max\!\left(10^{-30},\,\epsilon_{\chi}\max_k|\tilde{E}_k|\right),\\
j &\le j_{\mathrm{cut}}, \quad
j_{\mathrm{cut}}=\min\left\{j:\frac{\sum_{k=0}^j|\tilde{E}_k|^2}{\sum_{k}|\tilde{E}_k|^2}
\ge p_{\mathrm{cut}}\right\},\\
\omega_j &\ge \omega_{\min},\\
\omega_j &\le \omega_{\max},
\end{align}
with user controls $\epsilon_{\chi}$, $p_{\mathrm{cut}}$, and
$(\omega_{\min},\omega_{\max})=(E_{\min}^{(\chi)}e/\hbar,\,
E_{\max}^{(\chi)}e/\hbar)$. In the release workflow, the probe branch uses
$E_{\min}^{(\chi)}=1.5$ eV and $E_{\max}^{(\chi)}=3.0$ eV.

Optical observables are then derived from
$\tilde{n}(\omega)=\sqrt{1+\chi(\omega)}=n(\omega)+i\kappa(\omega)$:
\begin{align}
n(\omega)&=\mathrm{Re}\,\tilde{n}(\omega),\\
\kappa(\omega)&=\mathrm{Im}\,\tilde{n}(\omega),\\
\alpha(\omega)&=\frac{2\omega\kappa(\omega)}{c}.
\end{align}

\subsection{UPPE Propagation Operator Used in SBUP$^3$}
The temporal SBE fields are lifted to 2D using a Gaussian transverse envelope:
\begin{align}
E(x,t,0) &= \exp\!\left(-\frac{x^2}{w_0^2}\right)E_t(t),\\
P(x,t)   &= \exp\!\left(-\frac{x^2}{w_0^2}\right)P_t(t).
\end{align}
In $(k_x,\omega)$ space:
\begin{align}
k_0(\omega) &= \frac{n(|\omega|)\omega}{c},\\
k_z(k_x,\omega) &= \sqrt{\max\!\left(k_0^2-k_x^2,\,0\right)}.
\end{align}
For each propagation step $\Delta z$, the implemented update is
\begin{align}
\hat{E}^{(m+1)}
= \hat{E}^{(m)}e^{-ik_z\Delta z}
+ i\Delta z\,
\frac{\omega^2}{2k_z+\varepsilon}\,\hat{P},
\end{align}
with small denominator floor $\varepsilon=10^{-30}$. The polarization source is
treated as fixed along $z$ within a run, i.e., no intra-step material update is
performed inside UPPE. The propagated field is returned as
$E(x,t,z_{\mathrm{end}})=\mathrm{Re}\,\mathcal{F}^{-1}\{\hat{E}\}$.

\subsection{Self-Consistent Feedback Loop}
At coupling iteration $q$, UPPE output is sampled on axis:
\begin{align}
E_{\mathrm{out}}^{(q)}(t)=E^{(q)}(x\!\approx\!0,t).
\end{align}
Feedback is under-relaxed as
\begin{align}
E_{\mathrm{fb}}^{(1)}(t) &= E_{\mathrm{out}}^{(1)}(t),\\
E_{\mathrm{fb}}^{(q)}(t) &=
(1-\beta)E_{\mathrm{fb}}^{(q-1)}(t)+\beta E_{\mathrm{out}}^{(q)}(t), \quad q\ge2,
\end{align}
with $0<\beta\le1$.

Optional amplitude normalization then rescales $E_{\mathrm{fb}}^{(q)}$ to the
current SBE drive $E_{\mathrm{drv}}^{(q)}$:
\begin{align}
E_{\mathrm{fb}}^{(q)} \leftarrow
E_{\mathrm{fb}}^{(q)}
\frac{\mathrm{RMS}(E_{\mathrm{drv}}^{(q)})}{\mathrm{RMS}(E_{\mathrm{fb}}^{(q)})}
\quad \text{(RMS mode)},
\end{align}
or analogously with peak amplitudes in peak mode.

Convergence is checked by relative $L^2$ update norm,
\begin{align}
r_q=
\frac{\left\|E_{\mathrm{fb}}^{(q)}-E_{\mathrm{fb}}^{(q-1)}\right\|_2}
{\max\!\left(\left\|E_{\mathrm{fb}}^{(q-1)}\right\|_2,\epsilon_{\mathrm{abs}}\right)},
\end{align}
and iteration stops when $r_q<r_{\mathrm{tol}}$.

\subsection{Validation Metrics Used in This Work}
Two spectral moments are computed from normalized spectral power $S_j$:
\begin{align}
\mu &= \sum_j \hbar\omega_j S_j, \\
\sigma &= \sqrt{\sum_j (\hbar\omega_j-\mu)^2 S_j}.
\end{align}
The transverse energy diagnostic is
\begin{align}
\mathcal{E}(t)=\int |E(x,t)|^2\,dx,
\end{align}

reported through mean, standard deviation, and relative drift over the time
window. These diagnostics are used for release baselines, convergence sweeps,
and secondary cross-dataset checks under the same reporting protocol.

\section{Supplementary Figures}
\begin{figure}[t]
\centering
\includegraphics[width=0.5\textwidth]{ study_20260520/S1_postcal_fit_aspnes.pdf}
\caption{
Secondary post-calibration optical overlay (Aspnes): simulated probe-response
$|\alpha|$ and $n$ compared against the secondary validation reference over the
same analysis window used for primary validation.
}
\label{fig:supp_aspnes_overlay}
\end{figure}

\begin{figure}[t]
\centering
\includegraphics[width=0.5\textwidth]{ study_20260520/S2_probe_vs_drive_response.pdf}
\caption{
Drive-versus-probe susceptibility extraction comparison, including trusted
support masks.
}
\label{fig:supp_probe_drive}
\end{figure}

\begin{figure}[t]
\centering
\includegraphics[width=0.5\textwidth]{ study_20260520/S3_support_window_diagnostics.pdf}
\caption{
Support-window diagnostics from the release baseline, including probe trusted
support region and coverage summary.
}
\label{fig:supp_support_window}
\end{figure}



\section{DFT Provenance Tables}
Table~\ref{tab:appendix_dft_settings} summarizes the electronic-structure setup
used to generate DFT-informed parameters for SBUP$^3$. The workflow enforces
periodic non-$\Gamma$-only sampling for publication calibration and records both
ground-state and response artifacts. Table~\ref{tab:appendix_tddft_outputs}
lists legacy scalar quantities retained for compatibility.

Implementation-level DFT artifacts are organized as follows: ground-state
construction (\texttt{dft/gaas\_bulk\_gpaw.py}), dielectric-response extraction
(\texttt{dft/gaas\_dielectric\_response.py}), and transition reduction
(\texttt{dft/build\_transitions\_from\_dielectric.py}). Provenance for these
steps is recorded in \texttt{dft/gaas\_ground\_state\_meta.json},
\texttt{dft/gaas\_dielectric\_meta.json}, and
\texttt{dft/transitions\_meta.json}.


In the current implementation, the primary transition path is serialized in
\texttt{dft/transitions\_E\_eV.npy},
\texttt{dft/transitions\_f.npy}, and \texttt{dft/transitions\_meta.json}, then
consumed directly by \texttt{sbe/sbe\_rk4.py}. The legacy scalar files
\texttt{dft/E0\_eV.npy}, \texttt{dft/omega0\_rad\_s.npy},
\texttt{dft/f0.npy}, and \texttt{dft/d\_eff\_Cm.npy} remain available for the
compatibility path.



\clearpage
\bibliographystyle{apsrev4-2}
\bibliography{references}

\end{document}
